\documentclass[12pt]{article}
\usepackage{amsfonts,amsthm,amsmath,amssymb,mathtools}
\usepackage{mathrsfs}
\usepackage{enumitem}
\usepackage{tikz-cd}
\usepackage{geometry}
\usepackage{mlmodern}
\usepackage{xcolor}
\usepackage{pgfplots}

\geometry{letterpaper, portrait, margin=1in}
\pgfplotsset{compat=1.18}

\setcounter{section}{0}

\renewcommand{\thesection}{\S\arabic{section}}
\renewcommand{\thesubsection}{\arabic{section}}

\newtheorem{thm}{Theorem}
\newenvironment{theorem}[1]{
	\renewcommand{\thethm}{#1}
	\begin{thm}
		}{\end{thm}}

\newtheorem{lma}{Lemma}
\newenvironment{lemma}[1]{
	\renewcommand{\thelma}{#1}
	\begin{lma}
		}{\end{lma}}

\theoremstyle{definition}
\newtheorem{ex}{}
\newenvironment{exercise}[1]{
	\renewcommand{\theex}{#1}
	\begin{ex}
		}{\end{ex}}

\title{Homework Assignment 5}
\author{Connor Mooneyhan}
\date{September 26, 2025}

\begin{document}

\maketitle

\begin{ex}
	(2D.2)
	Prove that there exists a bounded set $A \subset \mathbb{R}$ such that $|F| \leq |A| - 1$ for every closed set $F \subset A$.
\end{ex}

\begin{ex}
	(2D.5)
	Prove that if $A \subset \mathbb{R}$ is Lebesgue measurable, then there exists an increasing sequence $F_1 \subset F_2 \subset \cdots$ of closed sets contained in $A$ such that \begin{equation*}
		\left| A \setminus \bigcup_{k=1}^\infty F_k \right| = 0.
	\end{equation*}
\end{ex}
\begin{proof}
	By Theorem 2.71, we can choose closed sets $G_1, G_2, \dots$ contained in $A$ such that $\left| A \setminus \bigcup_{k=1}^\infty G_k \right| = 0$.
	We want the sequence to be increasing, so define $F_n = \bigcup_{k=1}^n G_k$.
	Then $F_n \subset F_{n+1}$ for all $n$, and each $F_n$ is closed since it is a finite union of closed sets.
	Since $\bigcup_{k=1}^\infty F_k = \bigcup_{k=1}^\infty G_k$, we have \begin{equation*}
		\left| A \setminus \bigcup_{k=1}^\infty F_k \right| = \left| A \setminus \bigcup_{k=1}^\infty G_k \right| = 0.
	\end{equation*}
\end{proof}

\begin{ex}
	(2D.7)
	Prove that if $A \subset \mathbb{R}$ is Lebesgue measurable, then there exists a decreasing sequence $G_1 \supset G_2 \supset \cdots$ of open sets containing $A$ such that \begin{equation*}
		\left| \left( \bigcap_{k=1}^\infty G_k \right) \setminus A \right| = 0.
	\end{equation*}
\end{ex}
\begin{proof}
	By Theorem 2.71, there exist open sets $F_1, F_2, \dots$ containing $A$ such that \begin{equation*}
		\left| \left( \bigcap_{k=1}^\infty F_k \right) \setminus A \right| = 0.
	\end{equation*}
  Similar to the proof to 2D.5, we can now define $G_1, G_2, \dots$ by $G_n = \bigcap_{k=1}^n F_k$.
  Then each $G_n$ is open and contains $A$, and $\left\lbrace G_k \right\rbrace$ is a decreasing sequence, so \begin{equation*}
		\left| \left( \bigcap_{k=1}^\infty G_k \right) \setminus A \right| = \left| \left( \bigcap_{k=1}^\infty F_k \right) \setminus A \right| = 0.
  \end{equation*}
\end{proof}

\begin{ex}
	(2D.9)
	Prove that the collection of Lebesgue measurable subsets of $\mathbb{R}$ is dilation invariant.
	More precisely, prove that if $A \subset \mathbb{R}$ is Lebesgue measurable and $t \in \mathbb{R}$, then $tA$ (which is defined to be $\left\lbrace ta : a \in A \right\rbrace$) is Lebesgue measurable.
\end{ex}

\end{document}
